% Level 3, Week 2 -- Objects: Blueprints and Things
\lesson{2}{Objects: Blueprints and Things}{A class is a blueprint; objects are the things you build from it, and each one keeps its own data.}
\lessoninfo{Concepts: class, \_\_init\_\_, attributes, methods, self, dicts (review)}{Tools: Thonny, turtle, git}{CSTA: 3A-AP-17, 3A-AP-14}

\begin{objectives}
  \item The student can write a class with \texttt{\_\_init\_\_}, attributes and methods, and make several objects from it. (L3.G1)
  \item The student can explain what \texttt{self} means and fix a method that forgets it. (L3.G8)
  \item The student can use a list of objects to make turtle art in which every object keeps its own position, size and color. (L3.G3)
  \item The student can name game mechanics that are unkind to players and design a virtual pet without them. (L3.G9)
\end{objectives}

\begin{materials}
  \item Print: workbook chapter 2 and homework 2 (one per student).
  \item For the unplugged warm-up: 1 sheet of card, scissors, colored pencils, 4--6 sheets of plain paper.
  \item Before the session, cut a simple cookie shape (a star or gingerbread person) out of the card: that is your \enquote{cookie cutter}.
  \item On each Mac: Thonny with \texttt{grade-4-6/3-advanced/exercises/week02} open; Terminal in \texttt{Documents/git-quest} (\texttt{cd \textasciitilde/Documents/git-quest}).
\end{materials}

\section{At a glance}
\begin{timeplan}
  \timeslot{0--8}{Warm-up}{Unplugged \enquote{Cookie-cutter classes}: one cutter (class), many different cookies (objects).}
  \timeslot{8--18}{Concept}{From a dictionary to a class: attributes, methods, \texttt{\_\_init\_\_} and \texttt{self}.}
  \timeslot{18--28}{Guided demo}{Read the \texttt{Pet} class together, make two pets, and finish \texttt{play()} live.}
  \timeslot{28--52}{Hands-on build}{Pairs: Exercise 2.1 Pet Class, then 2.2 Bubble Art. Swap driver every 8 minutes. Bug hunt \enquote{Forgot self} on the way. Last 6 minutes: bubble free-build and a git commit.}
  \timeslot{52--60}{Wrap-up}{Show-and-tell of bubble pictures, kind-pet discussion, exit ticket, homework.}
\end{timeplan}

\section{Warm-up (0--8 min)}
\begin{unplugged}[Cookie-cutter classes]
Show the card cutter. Each student traces it twice on paper and cuts out two \enquote{cookies}. Then each student decorates their cookies differently (sprinkles, color, a face). On the back of each cookie they write three facts: \texttt{flavor}, \texttt{sprinkles} (a number) and \texttt{bites\_left = 3}.
Now play: \enquote{Ada's chocolate cookie, take a bite!} Only that cookie changes its \texttt{bites\_left} to 2. Ask:
\begin{itemize}[itemsep=1pt]
  \item \enquote{How many cutters do we have? How many cookies?} (One cutter, many cookies.)
  \item \enquote{When one cookie loses a bite, do the others change?} (No: each cookie has its own facts.)
\end{itemize}
\end{unplugged}

\section{Concept (8--18 min)}
\begin{teacherscript}
\enquote{In Level 2 you stored a game character in a dictionary:}
\begin{pycode}
pip = {"name": "Pip", "hunger": 5}
pip["hunger"] = pip["hunger"] - 3
\end{pycode}
\enquote{That works, but the actions, like feeding, live somewhere else in the program. A \textbf{class} keeps the data \emph{and} the actions together. The class is the cookie cutter. Each \textbf{object} is a cookie. The facts each object remembers are \textbf{attributes}. The actions it can do are \textbf{methods}: functions that live inside the class.}
\end{teacherscript}
Write this on paper or the board and point at each part as you explain it:
\begin{pycode}
class Pet:                          # the blueprint (cookie cutter)
    def __init__(self, name):       # runs once, when a pet is made
        self.name = name            # attribute: this pet's name
        self.hunger = 5             # attribute: this pet's hunger

    def feed(self):                 # method: an action
        self.hunger = self.hunger - 3

pip = Pet("Pip")                    # make an object (a cookie)
pip.feed()                          # Python runs feed with self = pip
\end{pycode}

\begin{analogy}[self means \enquote{my}]
When Pip runs \texttt{feed}, \texttt{self.hunger} means \enquote{\emph{my} hunger}. When Bo runs \texttt{feed}, the same code means Bo's hunger. \texttt{self} is how one blueprint can work for every object. Every method gets \texttt{self} as its first parameter, and Python fills it in for you: \texttt{pip.feed()} is really \texttt{Pet.feed(pip)}.
\end{analogy}
\texttt{\_\_init\_\_} has two underscores on each side (say \enquote{dunder init}). It is the setup method: it runs automatically when you write \texttt{Pet("Pip")}.

\section{Guided demo (18--28 min)}
Open \texttt{week02/ex-pet-class/starter.py} on your Mac with everyone watching.
\begin{workedexample}[The Pet class]
\begin{enumerate}[itemsep=2pt]
  \item Read \texttt{\_\_init\_\_} aloud. Ask: \enquote{Which four facts does every pet remember?} (name, species, hunger, happiness).
  \item Read \texttt{feed}. Ask: \enquote{Why do we check \texttt{if self.hunger < 0}?} (So the number never goes below 0.)
  \item Press Run and type \texttt{q} to quit the menu. Thonny remembers the class, so in the Shell (bottom of Thonny) try: \texttt{pip = Pet("Pip", "dragon")}, \texttt{pip.hunger}, \texttt{pip.feed()}, \texttt{pip.hunger}. Students predict each answer first.
  \item Live-code \texttt{play()} together, asking the students what to type at each line. Press Run, type \texttt{p}, and check that Pip's numbers change. The finished method (the full class is in the answer key):
\end{enumerate}
\pyfile[firstline=32,lastline=41]{../exercises/week02/ex-pet-class/solution.py}
\end{workedexample}

\section{Hands-on build (28--52 min)}
\begin{pairroles}[Driver and navigator]
The \textbf{driver} types. The \textbf{navigator} reads the README and the \texttt{TODO} comment aloud, and checks that every method has \texttt{self} as its first parameter. Swap at the 8-minute timer.
With \textbf{1 student}: you are the navigator and ask \enquote{whose hunger is that?} each time they type an attribute. With \textbf{3}: a \textbf{tester} runs \texttt{check\_my\_work.py} after each \texttt{TODO} and reads the hints; rotate roles. With \textbf{4}: two pairs.
\end{pairroles}
\begin{enumerate}[itemsep=2pt]
  \item \textbf{Exercise 2.1 Pet Class} (about 10 minutes): finish \texttt{play}, \texttt{mood} and \texttt{describe}. Run it and play with Pip; notice Bo never changes.
  \item \textbf{Exercise 2.2 Bubble Art} (about 10 minutes): finish the \texttt{Bubble} class and \texttt{make\_bubbles}. Forty bubble objects, each with its own \texttt{x}, \texttt{y}, \texttt{size} and \texttt{color}, fill the window.
\end{enumerate}

\begin{bughunt}[Forgot self]
Show this method and ask: \enquote{Find the two missing words.}
\begin{pycode}
    def play():
        happiness = happiness + 2
\end{pycode}
Python complains with \texttt{TypeError: Pet.play() takes 0 positional arguments but 1 was given}. That means \enquote{I tried to hand you the pet, but you have no \texttt{self} to catch it.} Fix: \texttt{def play(self):}. Then the next error, \texttt{UnboundLocalError}, means \texttt{happiness} is not the pet's: write \texttt{self.happiness = self.happiness + 2}.
\end{bughunt}

\begin{freebuild}[Bubble remix]
Last 6 minutes: change Bubble Art any way you like: more bubbles, a different background, bubbles that grow towards the top, or a new method such as \texttt{shine(t)} that adds a second highlight. Then commit it: \menu{File > Save as...} into \texttt{git-quest} with a new name such as \texttt{bubble\_art.py} (every week's file needs its own name, or it replaces the last one), then \texttt{git add bubble\_art.py} and \texttt{git commit -m "Add bubble art"}.
\end{freebuild}

\section{Wrap-up (52--60 min)}
Show-and-tell: each pair shows one bubble picture and names one attribute and one method.

\begin{safety}[Designing virtual pets kindly]
Ask: \enquote{Some pet apps say \enquote{Your pet is sad, you forgot it!} or let the pet get sick unless you pay or log in every day. How does that make players feel?} Talk about fair design: our pets never die, never blame the player, and never need money or daily streaks. Point to the comment at the top of \texttt{starter.py}: \enquote{Kind design: our pets never get sick, run away or make you feel guilty.} Programmers choose how their games treat people.
\end{safety}

\begin{formativecheck}[Exit ticket]
\textbf{Look for:} methods with \texttt{self} as the first parameter; attributes written as \texttt{self.something}.\\
\textbf{Ask:} \enquote{\texttt{pip = Pet("Pip", "dragon")}, \texttt{bo = Pet("Bo", "cat")}, then \texttt{pip.feed()}. What is \texttt{bo.hunger}? Why?}\\
\textbf{Answer:} 5. Feeding Pip changes only Pip's attributes; Bo is a different object.
\end{formativecheck}

\begin{reflection}[Talk or write]
\begin{itemize}[itemsep=1pt]
  \item What is the difference between a class and an object? Use your own example (not cookies!).
  \item Which object in a game you love could be a class? What attributes and methods would it have?
\end{itemize}
\end{reflection}

\section{Notes for the teacher}
\begin{commonmistakes}
  \item \textbf{Forgetting \texttt{self}} as the first parameter, or writing \texttt{hunger} instead of \texttt{self.hunger}. Ask \enquote{whose hunger?}
  \item \textbf{Typing \texttt{\_\_int\_\_} or \texttt{\_init\_}} instead of \texttt{\_\_init\_\_}. The error is \texttt{TypeError: Pet() takes no arguments}. Count the underscores: two on each side.
  \item \textbf{Calling a method without parentheses}: \texttt{pip.feed} does nothing. It must be \texttt{pip.feed()}.
  \item \textbf{Methods not indented} under the class, so they become ordinary functions. All methods line up 4 spaces in.
  \item \textbf{Thinking \texttt{Pet} is the pet.} \texttt{Pet} is the blueprint; \texttt{pip} is the pet.
\end{commonmistakes}

\begin{differentiation}
\supportpath Do Exercise 2.1 only, one \texttt{TODO} at a time, running \texttt{check\_my\_work.py} after each. Keep the labeled class picture from the workbook beside the keyboard.
\extensionpath Add a \texttt{nap()} method and a third pet, or give \texttt{Bubble} a method \texttt{is\_touching(other)} that returns \texttt{True} when two bubbles overlap (distance between centers less than the sum of their radii).
\end{differentiation}

\begin{stretchchallenge}[Floating bubbles]
Use \texttt{screen.ontimer} (Level 2) to animate the bubbles: every 50 milliseconds, call \texttt{move(0, 2)} on each bubble, clear the pen, and draw them all again. Bubbles that float off the top start again at the bottom.
\end{stretchchallenge}

\begin{teachernote}[If you are short on time]
Keep the cookie warm-up short (one cookie each). Do the Pet class fully; Bubble Art can be finished at home or next week. The bug hunt can be done in the Shell in 2 minutes during the Pet exercise.
\end{teachernote}
